\documentclass[10pt]{beamer}
% !TEX program = xelatex
\usetheme{metropolis}
\usepackage{appendixnumberbeamer}
\usepackage{booktabs}
\usepackage[scale=2]{ccicons}
\usepackage{pgfplots}
\usepgfplotslibrary{dateplot}
\usepackage{xspace}
\newcommand{\themename}{\textbf{\textsc{metropolis}}\xspace}
\usepackage{caption}
\usepackage{fontspec}
\usepackage{xeCJK}
\setCJKmainfont{Noto Serif CJK SC}
\setCJKsansfont{Noto Sans CJK SC}


% --- Custom Header for UCSD Tailoring ---
\title{Behavioral Discounting of Delayed Protection: Evidence from a
Pandemic Vaccination Discrete Choice Experiment}

\subtitle{Job Market Paper --- Postdoctoral Candidate Presentation \\
\vspace{0.2em}
\small Herbert Wertheim School of Public Health | UC San Diego}

\author{Yichao Jin}
\institute{The University of Texas at Dallas}
\date{\today}

\begin{document}

\maketitle

% -----------------------------------------------------------------------------
% SECTION I: INTRODUCTION
% -----------------------------------------------------------------------------
\section{Introduction}

\begin{frame}{Motivation: The Intertemporal Dimension of Health}
    \begin{itemize}
        \item Global vaccination research: Focused on "Hesitancy" vs. "Acceptance."
        \item \textbf{Our Contribution}: Focus on the \alert{Intertemporal Dimension}.
        \item Does the \textit{vaccine delay} to protection significantly discount the perceived value of vaccination?
        \item This is critical for logistics, booster scheduling, and managing future outbreaks.
    \end{itemize}
\end{frame}

% --- 第 1 页：提出问题 ---
\begin{frame}{Research Questions}
  \begin{itemize}
    \item \textbf{Identification}: Recover structural parameters ($\delta, \kappa$) for delayed protection.
    \item \textbf{Mechanisms}: Distinguish Exponential vs. \alert{Hyperbolic} discounting.
    \item \textbf{Heterogeneity}: Map preference variation across trust and income subgroups.
  \end{itemize}
\end{frame}

% --- 第 2 页：给出答案 ---
\begin{frame}{Preview: Core Take-home Messages}
  \begin{enumerate}
    \item \textbf{Mechanism}: Vaccination choices are best captured by \alert{Hyperbolic Discounting} ($\kappa > 0$).
    \item \textbf{Psychological Cost}: The "Behavioral Delay Multiplier" (\textbf{BDM}) is \alert{1.29} --- waiting costs 29\% more than rational ($\kappa=0$) expectations predict.
    \item \textbf{Welfare Impact}: The shadow price of waiting (\textbf{MWTA}) is approx. \alert{111 RMB/month} ($\approx$\$15).
  \end{enumerate}
  \vfill
  \begin{center}
    \textit{Trust and income significantly moderate these intertemporal costs, with high-trust individuals exhibiting lower psychological discounting.}
  \end{center}
\end{frame}
% -----------------------------------------------------------------------------
% SECTION II: METHODOLOGY
% -----------------------------------------------------------------------------
% -----------------------------------------------------------------------------
% SECTION II: METHODOLOGY: THE BEHAVIORAL TOOLBOX
% -----------------------------------------------------------------------------
\section{Methodology: The Behavioral Toolbox}

\begin{frame}{Target Population and Study Setting}
\begin{itemize}
\item \textbf{The Epicenter Context (Wuhan)}:
\begin{itemize}
\item Unique historical proximity to the initial COVID-19 outbreak.
\item High \alert{risk-salience} among residents provides a robust environment for eliciting intertemporal preferences.
\end{itemize}
\end{itemize}
\end{frame}

\begin{frame}{Survey Design and Structure}
\begin{itemize}
\item \textbf{Platform}: Online survey panel in Wuhan ($N=1,027$).
\item \textbf{Instrument Structure}:
\begin{enumerate}
\item \textbf{Socio-Demographics}: Age, gender, income, education, residence.
\item \textbf{Health \& Attitudes}: Biological vulnerability, institutional trust, preventive-health routines.
\item \textbf{The DCE}: 6 choice tasks per respondent (drawn from a D-efficient design).
\item \textbf{Quality Control}: Attention filters, time-to-complete, and dominance checks.
\end{enumerate}
\end{itemize}
\end{frame}


\begin{frame}{Framing Strategy: COVID-like Diseases (CLD)}
  \begin{itemize}
    \item \textbf{Scenario Design}: Hypothetical respiratory disease with pandemic potential.
    \item \textbf{Why CLD? (The Rationale)}:
      \begin{itemize}
        \item \alert{Minimizes Anchoring}: Decouples preferences from specific COVID-19 controversies.
        \item \alert{Neutralizes Stigma}: Reduces social desirability bias in behavior reporting.
      \end{itemize}
    \item \textbf{External Validity}:
      \begin{itemize}
        \item Recovers structural parameters for \alert{Pandemic Preparedness}.
        \item Generalizable to other infectious diseases (e.g., flu, future outbreaks).
      \end{itemize}
  \end{itemize}
\end{frame}

\begin{frame}{Attribute and Level Development}
\begin{itemize}
\item \textbf{Stage 1: Evidence Synthesis}: Review of 50+ papers and policy reports on vaccine hesitancy and time preferences.
\item \textbf{Stage 2: Expert Consultation}: Qualitative refinement with public health practitioners to ensure policy relevance.
\item \textbf{Stage 3: Pilot Study ($n=60$)}:
\begin{itemize}
\item Cognitive interviews to verify comprehension of "Time to Protection."
\item Used pilot data to elicit \alert{priors} for the final experimental design.
\end{itemize}
\end{itemize}
\end{frame}

\begin{frame}{Experimental Design: Attributes and Levels}
\begin{table}
\centering
\small
\begin{tabular}{ll}
\toprule
Attribute & Levels \\
\midrule
\textbf{Vaccination Delay} ($T$) & \alert{0, 1, 3, 6 months} \\
Vaccine Efficacy ($E$) & 50\%, 70\%, 95\% \\
Side Effects ($S$) & Mild, Moderate, Severe \\
Vaccine Origin ($O$) & Domestic, Imported \\
Incentives ($M$) & 50, 200, 800 RMB \\
\bottomrule
\end{tabular}
\caption{Final attributes justified by policy roll-out constraints.}
\end{table}
\end{frame}

\begin{frame}{DCE Construction: Sample Choice Task}
\begin{table}
\centering
\small
\begin{tabular}{|l|c|c|c|}
\hline
\textbf{Attribute} & \textbf{Vaccine A} & \textbf{Vaccine B} & \textbf{Opt-out} \\
\hline
Vaccine Delay & 0 months & 3 months & \\
Efficacy & 95\% & 95\% & \\
Side Effects & Mild & Mild & \textbf{I would} \\
Origin & Domestic & Imported & \textbf{choose} \\
Incentive & 50 RMB & 200 RMB & \textbf{neither.} \\
\hline
\textbf{Your Choice?} & $\square$ & $\square$ & $\square$ \\
\hline
\end{tabular}
\end{table}
\begin{itemize}
\item \textbf{The Task}: Choose the preferred vaccine profile or opt-out.
\item \textbf{Design Efficiency}: D-efficient design using pilot priors.
\end{itemize}
\end{frame}



\begin{frame}{Theoretical Foundation: Random Utility Framework}
    \begin{itemize}
        \item \textbf{The Core Assumption}: Individual $n$ chooses the vaccine that maximizes:
    \end{itemize}
    
    \begin{equation}
        U_{njt} = \underbrace{\text{Discounting}(T_j; \delta, \kappa)}_{\text{Behavioral Core}} \cdot \underbrace{V(X_j, M_j)}_{\text{Attributes}} + \epsilon_{njt}
    \end{equation}
    
    \vspace{0.4cm}
    \begin{itemize}
        \item \textbf{Structural Specificity}:
        \begin{itemize}
            \item \textbf{Exponential}: $\text{Discounting} = \delta^T$
            \item \textbf{Hyperbolic}: $\text{Discounting} = (1 + \kappa T)^{-1}$
        \end{itemize}
        \item \textbf{The Error Term ($\epsilon$)}: IID Extreme Value Type I (leads to Logit probabilities).
    \end{itemize}
    \vfill
    \centering \textit{$\implies$ We estimate this via MXL to allow for preference heterogeneity.}
\end{frame}
% SECTION III: RESULTS (JMP)
% -----------------------------------------------------------------------------
% --- Section III: Results I: Completed Study (JMP) ---
% --- SECTION III: RESULTS I (THE BEHAVIORAL DISCOVERY) ---
\section{Results I: Behavioral Discounting Findings}

\begin{frame}{Result 1: Hyperbolic Model Best Explains Choice Data}
    \begin{itemize}
        \item \textbf{The Contest}: We tested competing intertemporal functional forms.
    \end{itemize}
    \begin{table}[]
        \centering \small
        \begin{tabular}{lccc}
            \hline
            \textbf{Model} & \textbf{Log-Likelihood} & \textbf{AIC} & \textbf{Evidence} \\
            \hline
            Exponential (Standard) & -4,521.2 & 9,046.4 & -- \\
            \textbf{Hyperbolic (Behavioral)} & \textbf{-4,288.5} & \textbf{8,581.0} & \alert{\textbf{Best Fit}} \\
            \hline
        \end{tabular}
    \end{table}
    \vfill
    \begin{exampleblock}{\centering \textbf{Conclusion: Present-Bias is Real}}
        \centering \small Human vaccine timing is not a constant-rate process; it is dominated by a \alert{disproportionate penalty for immediate delays}.
    \end{exampleblock}
\end{frame}

\begin{frame}{The Mechanism: A Sharp "Present-Bias" Penalty ($\kappa$)}
    \begin{columns}[T]
        \begin{column}{0.48\textwidth}
            {\large \textbf{Estimated Parameter}}
            \begin{exampleblock}{\centering $\hat{\kappa} = \mathbf{0.225}$ ($p<0.01$)}
            \end{exampleblock}
            \vspace{0.5cm}
            \textbf{The Behavioral Hurdle}:
            \begin{itemize}
                \item Captures the \alert{sharp drop} in value the second a wait starts.
                \item Explains why even short logistics delays kill acceptance.
            \end{itemize}
        \end{column}
        
        \begin{column}{0.52\textwidth}
            \centering
            \setlength{\unitlength}{0.9cm}
            \begin{picture}(5,4)
                \put(0,0){\vector(1,0){4.5}} \put(4.6,-0.1){$t$}
                \put(0,0){\vector(0,1){3.5}} \put(-0.8,3.3){Value}
                \color{red} \thicklines
                \qbezier(0,3)(0.2,0.8)(4,0.5)
                \color{black} \thinlines
                \put(0.1,2.8){\vector(1,-2){0.4}}
                \put(0.6,2.2){\scriptsize \alert{Instant Drop}}
            \end{picture}
            \vfill
            \small \textbf{Takeaway}: It's the \textbf{structural pain} of "not having it now."
        \end{column}
    \end{columns}
\end{frame}


\begin{frame}{The Behavioral Delay Multiplier (BDM)}
    \begin{columns}[T]
        \begin{column}{0.48\textwidth}
            \begin{block}{Concept}
                Rescaling $\hat{\kappa}$ into a transparent policy metric:
                \centering \large $BDM = \frac{1}{1 - \hat{\kappa}}$
            \end{block}
            \vspace{0.3cm}
            \begin{exampleblock}{Interpretation}
                \small A \textbf{14-day} wait feels like an \alert{\textbf{18.1-day}} delay in perceived utility.
            \end{exampleblock}
        \end{column}
        
        \begin{column}{0.52\textwidth}
            \centering
            \begin{beamercolorbox}[sep=1.5em, wd=\linewidth, center]{block title}
                \Huge \alert{\textbf{BDM = 1.29}}
            \end{beamercolorbox}
            \vspace{0.3cm}
            \small \textbf{The "Psychological Markup"} \\
            \textit{Waiting costs 29\% more than rational expectations predict.}
        \end{column}
    \end{columns}
    \vfill
    \centering \textit{This 29\% markup explains why even high-efficacy vaccines face rollout resistance.}
\end{frame}

\begin{frame}{Deep Heterogeneity in Time Preferences}
    \begin{itemize}
        \item \textbf{MXL Estimation}: Significant Standard Deviations ($p<0.01$) for $\delta$ and $\kappa$, confirming \alert{latent preference variation}.
    \end{itemize}
    \begin{table}[]
        \centering \small
        \begin{tabular}{lcc}
            \toprule
            Subgroup Characteristics & Impact on $\kappa$ (Present-Bias) & Policy Barrier \\
            \midrule
            \textbf{Low Institutional Trust} & \alert{+55\% Inflation} & High Behavioral Hurdle \\
            \textbf{Chronic Medical Condition} & \alert{+23\% Sensitivity} & Biological Urgency \\
            Urban Residence & +29\% Sensitivity & Lifestyle Friction \\
            \bottomrule
        \end{tabular}
    \end{table}
    \vfill
    \begin{exampleblock}{The "Equity Gap"}
        \small Structural preferences are stratified by \alert{Trust} and \alert{Vulnerability}. A fixed-delay policy acts as a "regressive tax" on high-risk and low-trust populations.
    \end{exampleblock}
\end{frame}










% --- DATA QUALITY & VALIDITY (SHI LAB EDITION) ---

\begin{frame}{Data Quality: The Evidence is Clean}
\vspace{0.6cm}
\begin{columns}[T]
\begin{column}{0.48\textwidth}
\begin{beamercolorbox}[sep=0.8em, wd=\linewidth, center]{block title}
\Huge \textbf{92\%} \\
\small \textbf{DCE Dominance Pass}
\end{beamercolorbox}
\vspace{0.4cm}
\centering \small $\implies$ High Rationality
\end{column}

\begin{column}{0.48\textwidth}
\begin{beamercolorbox}[sep=0.8em, wd=\linewidth, center]{block title}
\Huge \textbf{85.5\%} \\
\small \textbf{Test-Retest Reliability}
\end{beamercolorbox}
\vspace{0.4cm}
\centering \small $\implies$ Reliable Elicitation
\end{column}
\end{columns}
\vfill
\begin{itemize}
\item \textbf{Result}: Strong internal validity ensures these aren't random choices.
\end{itemize}
\end{frame}

\begin{frame}{Internal Validity: Consistent with Behavioral Theory}
    \begin{center}
        \includegraphics[width=0.65\textwidth]{correlation among Individual variables4.png}
    \end{center}
    \vfill
    \begin{exampleblock}{\centering \textbf{The Core Signal}}
        \centering \small Structural parameters ($\kappa, \delta$) significantly correlate with \alert{Trust} and \alert{Vulnerability}, proving preferences are linked to real-world traits.
    \end{exampleblock}
\end{frame}





% -----------------------------------------------------------------------------
% SECTION IV: RESULTS II: MONETARY VALUATION & POLICY EVALUATION
% -----------------------------------------------------------------------------
\section{Results II: Monetary Valuation \& Policy Evaluation}

\begin{frame}{Essay 2 Methodology: The Monetary Bridge}
    \begin{itemize}
        \item \textbf{Objective}: Translating abstract parameters ($\delta, \kappa$) into \alert{RMB Units}.
        \item \textbf{Approach}: Structural estimation of the "Shadow Price of Waiting."
    \end{itemize}
    
    \begin{exampleblock}{The Estimand: Marginal Willingness to Accept (MWTA)}
        \centering
        \large $MWTA = - \frac{\partial U / \partial \text{Delay}}{\partial U / \partial \text{Incentive}} = - \frac{\beta_{delay}}{\beta_{incentive}}$
    \end{exampleblock}
    
    \begin{itemize}
        \item \textbf{Structural Identification}: Identified via experimental trade-offs between waiting time (0--6 mo) and cash (0--800 RMB).
        \item \textbf{Why it matters}: Converts behavioral friction into a \alert{budget-ready metric} for public health planners.
    \end{itemize}
\end{frame}

\begin{frame}{Essay 2 Results: Heterogeneous Shadow Prices}
    \begin{columns}
        \begin{column}{0.55\textwidth}
            \centering
            \includegraphics[width=\textwidth]{Subgroup differences in time sensitivity.png}
            \captionof{figure}{\footnotesize{MWTA across subgroups}}
        \end{column}
        \begin{column}{0.45\textwidth}
            \begin{itemize}
                \item \textbf{Average Shadow Price}: $\approx$ \textbf{111 RMB/month}.
                \item \textbf{The Trust Premium}: 
                \begin{itemize}
                    \item Low-trust individuals require \alert{2.2x higher} compensation to tolerate the same delay.
                \end{itemize}
                \item \textbf{Policy Insight}: Trust acts as a "behavioral stabilizer" that reduces the financial cost of rollout frictions.
            \end{itemize}
        \end{column}
    \end{columns}
\end{frame}

\begin{frame}{Simulation: Reducing Delay vs. Increasing Subsidies}
\begin{columns}[T]
% 左侧图
\begin{column}{0.6\textwidth}
\centering
\vspace{0.6cm}
\includegraphics[width=\textwidth]{essay3_policy_tradeoff_figure (1).png}
\end{column}

% 右侧内容
\begin{column}{0.4\textwidth}
\vspace{0.8cm}
\textbf{Simulation Results:}
\begin{itemize}
\item \textbf{Efficiency}: Reducing delay yields larger welfare gains than cash alone.
\item \textbf{Equity}: Targeted subsidies are needed for low-trust groups.
\end{itemize}
\vspace{0.6cm}
\begin{quote}
\color{orange}
\textbf{Key Insight:} Shortening the "time tax" is the most effective policy lever.
\end{quote}
\end{column}
\end{columns}
\end{frame}

% -----------------------------------------------------------------------------
% SECTION V: BROADER RESEARCH & FUTURE AGENDA
% -----------------------------------------------------------------------------
\section{Broader Research \& Future Agenda}

\begin{frame}{Research Portfolio: Beyond the Dissertation}
    \begin{itemize}
        \item \textbf{Completed Studies (Dissertation Overview)}
            \begin{itemize}
                \item \textbf{Essay 1}: Structural recovery of health discounting parameters ($\delta, \kappa$).
                \item \textbf{Essay 2}: Monetary valuation (MWTA) and "Shadow Price of Waiting."
                \item \textbf{Essay 3}: Equity-sensitive simulations for policy optimization.
            \end{itemize}
        \item \textbf{Methodological Strengths (Your Future Lab's Assets)}
            \begin{itemize}
                \item \textbf{Structural Econometrics}: FIML/Simulated Likelihood; Latent Class modeling.
                \item \textbf{Experimental Design}: D-efficient DCEs; Incentive-compatible mechanisms.
                \item \textbf{Computational Mastery}: R (Expert), Stata, Python; Large-scale data scraping.
            \end{itemize}
    \end{itemize}
\end{frame}

\begin{frame}{Future Project 1: Addiction \& Intertemporal Choice}
    \begin{itemize}
        \item \textbf{Objective}: Modeling the "Relapse Loop" in substance use cessation.
        \item \textbf{Structural Framework}:
            \begin{itemize}
                \item Applying \alert{Hyperbolic Discounting} models to longitudinal tobacco/cannabis data.
                \item Identifying the structural point where craving utility overrides long-term health utility.
            \end{itemize}
        \item \textbf{Policy Goal: Dynamic Incentive Design}
            \begin{itemize}
                \item Testing \alert{Front-loaded Rewards} for the critical "Day 1--7" window.
                \item Using BDM-based mechanisms for user pre-commitment contracts.
            \end{itemize}
    \end{itemize}
\end{frame}

\begin{frame}{Future Project 2: Digital Health Equity}
    \begin{itemize}
        \item \textbf{Objective}: Quantifying the regressive \alert{"Cognitive Time-Tax"} of digital nudges.
        \item \textbf{Methodology: Experimental Elicitation of Nudge Fatigue}
            \begin{itemize}
                \item Using DCEs to estimate the \alert{MWTA for cognitive load} among Medicaid populations.
                \item Attributes: Frequency, timing, and framing (loss vs. gain).
            \end{itemize}
        \item \textbf{Impact}:
            \begin{itemize}
                \item Designing "nudge-neutral" interventions for high-risk populations.
                \item Reducing administrative/behavioral burden in preventive care.
            \end{itemize}
    \end{itemize}
\end{frame}

\begin{frame}{Future Alignment: The "Shi Lab" Synergy}
    \begin{itemize}
        \item \textbf{Substance Use Policy}: Applying my structural toolkit to UCSD's world-class tobacco and cannabis research.
        \item \textbf{Health Equity}: Bridging behavioral modeling with the Herbert Wertheim School’s mission for underserved populations.
    \end{itemize}
\end{frame}

\section{Conclusion}

\begin{frame}{Conclusion: A Behavioral Bridge to Policy}
    \begin{itemize}
        \item \textbf{Main Contributions}:
            \begin{itemize}
                \item Structural evidence of \alert{Hyperbolic Discounting} in vaccination timing.
                \item Introduction of policy-ready metrics: \alert{BDM} and \alert{MWTA}.
                \item Demonstrated the crucial role of \alert{Institutional Trust} as a behavioral stabilizer.
            \end{itemize}
        \item \textbf{Strategic Fit}:
            \begin{itemize}
                \item Ready to leverage this structural modeling toolkit to support the \textbf{Shi Lab’s} high-impact missions in substance use and health policy.
                \item Committed to the \textbf{UCSD Herbert Wertheim School’s} pursuit of evidence-based health equity.
            \end{itemize}
    \end{itemize}
    \vfill
    \begin{center}
        \small \textit{Thank you for your time. I look forward to your questions.}
    \end{center}
\end{frame}

\begin{frame}[standout]
    \centering
    {\Huge Thank you!}\\
    \vspace{1.5em}
    \textbf{Yichao Jin}\\
    \small Yichao.Jin@UTDallas.edu
\end{frame}

\end{document}
